\documentclass[../../../include/open-logic-section]{subfiles}

\begin{document}

\olfileid{sth}{spine}{zf}
\olsection{$\Z$ تے $\ZF$۔ اک اہم پڑاؤ}

بنیاد دے مسلّمے نال اسیں اک ہور اہم پڑاؤ اُتے پہنچدے آں۔ اسیں
$\Zminus$ تے $\ZFminus$ نظریے ویکھے نیں تے آکھیا سی کہ ایہ کجھ
نظریے نیں جیہناں وچوں اک خاص شے ’’گھٹائی‘‘ اے۔ اوہ خاص شے بنیاد دا
مسلّمہ اے۔ سو:

\begin{defn}
نظریہ $\Z$، $\Zminus$ وچ بنیاد دا مسلّمہ شامل کردا اے۔ سو اوہدے
مسلّمے ایہ نیں۔ عنصراں نال تعیّن، اتحاد، جوڑیاں، ذیلی سیٹاں دے سیٹ،
لامتناہیت، بنیاد، تے وکھ کرن دے خاکے دیاں ساریاں صورتاں۔

نظریہ $\ZF$، $\ZFminus$ وچ بنیاد دا مسلّمہ شامل کردا اے۔
دوجے لفظاں وچ، $\ZF$،~$\Z$ وچ بدل رکھن دے خاکے دیاں ساریاں صورتاں
شامل کردا اے۔
\end{defn}

فیر وی اک سوال تہانوں سُجھیا ہو سکدا اے۔ جے باقاعدگی، $\ZFminus$ دے
لحاظ نال، بنیاد دے ہم ارز اے تے باقاعدگی دا جواز صاف اے، تاں
اینی لمّی راہ کیوں لئے، تے باقاعدگی دی تھاں بنیاد نوں اپنا بنیادی مسلّمہ کیوں منّیئے؟

تاریخی وجوہاں (کس نے کیہ تے کدوں بیان کیتا) اک پاسے رکھ کے، بنیادی وجہ
ایہ اے کہ بنیاد دا مسلّمہ~$V_\alpha$ والے سیٹاں دی تعریف ورتے بناں پیش
ہو سکدا اے۔ اوہ تعریف \olref[recursion]{sec} دے سارے کم اُتے ٹکی سی۔
اوہدا جواز وکھاؤن لئی سانوں لامتناہی درجیاں والی بازگشت ثابت کرنی پئی سی۔
پر اوہ بازگشت ثابت کرن لئی اسیں \emph{بدل رکھن دا اصول} ورتیا سی۔
سو بھانویں بنیاد تے باقاعدگی $\ZFminus$ دے لحاظ نال ہم ارز نیں،
$\Zminus$ دے لحاظ نال ہم ارز نہیں۔

اصل وچ گل ایس سادی ٹِپّن توں وی اگے جاندی اے۔ بھانویں ایہ کتاب دے
دائرے توں کافی باہر اے، پتا لگدا اے کہ $\Zminus$ تے $\Z$ دونیں
$V_\alpha$ والے سیٹاں دی تعریف لئی بہتے کمزور نیں۔ سو جے تسی صرف
$\Z$ وچ کم کر رہے او، تاں باقاعدگی (جیویں اسیں بیان کیتی اے) دا
\emph{مطلب} ای نہیں بندا۔ ایس لئی ساڈا مقرر کیتا مسلّمہ بنیاد اے، باقاعدگی نہیں۔

ہُن توں اسیں $\ZF$ وچ کم کراں گے (جے کجھ ہور نہ آکھیا جاوے)، تے
ایس بارے ہور کوئی ٹِپّن نہیں کراں گے۔

\end{document}
